\section{Introduction}
% To-do: Establish that FSDP is acronym for Fully Sharded Data Parallel
% We generally refer to FSDP to describe the concept as a whole. We specify FullyShardedDataParallel if referring to the Python class itself.

The magnitude of neural network models is growing at an unprecedented rate, facilitating breakthroughs across a wide spectrum of domains. Upon inception, the 175-billion-parameter GPT-3~\cite{gpt3} model set a new record for almost all Natural Language Processing tasks. The product applications constructed on top of GPT models~\cite{chatgpt} have quickly demonstrated their potential to revolutionize the entire industry. Modern large scale recommendation models~\cite{mudigere2021high, DHEN} can reach beyond 1 trillion parameters, replete with rapidly growing dense layer components. These models power applications that serve multi-billions of users every day. As large neural networks continue to push the limits of science and technology, an industry-grade tool to simplify the training of such models with high efficiency would help expedite the progress.

In recent years, the community has introduced and investigated numerous advanced methodologies to enlarge neural network models. Pipeline parallelism~\cite{gpipe, kim2020torchgpipe, pipetransformer, li2021terapipe, narayanan2019pipedream} partitions a model instance into stages and distributes stages across multiple devices, where activations and gradients are communicated across stage boundaries. Tensor parallelism~\cite{flexflow, narayanan2021efficient, gspmd, yuan2021oneflow} shards model parameters, conducts partial computation on individual devices and communicates activations at required layer boundaries. Zero-Redundancy parallelism~\cite{zero, zero:offload, xu2020automatic} shards parameters as well but communicates parameters on-demand to recover their unsharded form and executes the model as if it were replicated on every device. The aforementioned techniques have served as the fundamental building blocks to enable the training of large neural networks across various applications. Nevertheless, two challenges still persist. Firstly, some of these methods are tightly integrated with specific model architectures, which hinder them from being utilized as a generic solution for training large models. Secondly, some of these techniques are built on top of rapidly-evolving internal interfaces of underlying machine learning frameworks, which become vulnerable to changes in framework implementations. Therefore, it is more robust and efficient to have a native solution co-designed with the core functionalities of machine learning frameworks. Additionally, constructing such a solution in a composable and customizable manner could potentially facilitate the community's future innovations as well.

This paper presents PyTorch~\cite{pytorch:nips:2019} Fully Sharded Data Parallel (FSDP), which enables the training of large-scale models by sharding model parameters. The FSDP algorithm is motivated by the ZeroRedundancyOptimizer~\cite{zero, zero:offload} technique from DeepSpeed but with a revised design and implementation that is aligned with the other components of PyTorch. FSDP breaks down a model instance into smaller units and then flattens and shards all of the parameters within each unit. The sharded parameters are communicated and recovered on-demand before computations, and then they are immediately discarded afterwards. This approach ensures that FSDP only needs to materialize parameters from one unit at a time, which significantly reduces peak memory consumption. The design and implementation of FSDP faces the following challenges.

\begin{itemize}
    \item \textbf{User Experience} is critical for achieving broad adoption. When working on prior PyTorch distributed training features such as \code{DistributeDataParallel} (DDP)~\cite{ddp}, we observed that aligning the user experience of distributed training with that of local training can significantly lower the learning barrier. Techniques like DDP require the model to be replicated on every device, which implies that the entire model can be constructed on the target device. However, although FSDP can easily adopt DDP's API design, large models might not fit into one GPU device and therefore cannot even be initialized efficiently. 
    \item \textbf{Hardware Heterogeneity} often exists in modern GPU clusters, whereby interconnects are partitioned into high-bandwidth islands within each machine and low-bandwidth mesh across machines. Additionally, there may be further hierarchical structures at the rack or pod levels. Consequently, the design of FSDP must accommodate such heterogeneity and optimize accordingly. 
    \item \textbf{Resource Utilization} is usually tightly linked with capital and operational expenditures, especially for companies that depend on large GPU clusters to power their mission-critical systems. To ensure that GPU devices remain fully utilized during distributed training, it is essential to minimize downtime caused by non-computational operations.
    \item \textbf{Memory Planning} plays a crucial role in large model training. PyTorch makes GPU memory block allocation efficient and transparent through caching. Frequent memory defragmentations can significantly slow down training, which becomes particularly acute when working with large models. In such scenarios, practitioners typically seek to saturate GPU memory as much as possible to accommodate the largest batches or models. However, operating near GPU memory capacity significantly increases the chance to trigger defragmentations.
\end{itemize}

FSDP tackles the aforementioned challenges through a variety of techniques. Firstly, to improve user experience, FSDP introduces deferred initialization that allows users to create a model instance on a dummy device and record operations invoked during initialization. Then, the model can be initialized and sharded unit by unit by replaying the recorded operations on a real GPU device. With this technique, FSDP can provide similar user experiences as local training, while effectively scaling large models. Secondly, FSDP offers configurable sharding strategies that can be customized to match the physical interconnect topology of the cluster to handle hardware heterogeneity. Thirdly, although parameter sharding design inevitably inserts communications, which might block computations and introduces bubbles during execution, FSDP can squeeze out bubbles using an abundant set of tools to aggressively overlap communication with computation through operation reordering and parameter prefetching. Lastly, FSDP optimizes memory usage by prudently restricting the amount of blocks allocated for inflight unsharded parameters and suspending CPU execution if necessary.

We evaluated the performance of FSDP on various models including popular language models and recommendation system models, utilizing up to 512 80GB A100 GPUs. The experiments showed that FSDP can achieve similar performance to that of DDP on small models. Beyond that FDSP can facilitate significantly larger models with near-linear scalability in terms of TFLOPS. FSDP is currently a beta feature as of PyTorch 2.0 release, and has been battle-tested by both industrial and research applications. 

To simplify presentation, the rest of this paper uses FSDP to refer to the techniques in general and \code{FullyShardedDataParallel} to denote the Python implementation. The remainder of the paper is organized as follows. Section~\ref{sec:background} introduces background on some popular distributed training techniques. Section~\ref{sec:design} and Section~\ref{sec:implementation} elaborate system design and implementation details. Evaluations are presented in Section~\ref{sec:evaluation}. Section~\ref{sec:related} surveys related work, and Section~\ref{sec:discussion} discusses topics related to FSDP but falls outside of FSDP core. Finally, Section~\ref{sec:conclusion} concludes the paper. 




